% Generated draft: requires mathematical and source review.
\documentclass[11pt]{article}
\usepackage[margin=25mm]{geometry}
\usepackage{fontspec}
\setmainfont{Latin Modern Roman}
\usepackage{amsmath,amssymb,mathtools,unicode-math}
\setmathfont{Latin Modern Math}
\usepackage{enumitem,needspace,xurl,hyperref}
\setlength{\parindent}{0pt}
\setlength{\parskip}{5pt}
\setlength{\emergencystretch}{3em}
\newcommand{\N}{\mathbb{N}}
\newcommand{\Z}{\mathbb{Z}}
\newcommand{\Q}{\mathbb{Q}}
\newcommand{\R}{\mathbb{R}}
\newcommand{\C}{\mathbb{C}}
\newcommand{\F}{\mathbb{F}}
\newcommand{\A}{\mathbb{A}}
\newcommand{\PP}{\mathbb{P}}
\newcommand{\E}{\mathbb{E}}
\newcommand{\eps}{\varepsilon}
\newcommand{\dd}{\,\mathrm{d}}
\begin{document}
\section*{UnsolvedMath ID: 701383}
\textbf{Source problem code:} catalogue entry
\subsection{Definitions and mathematical statement}
Let \(G\) be a free group, with no restriction on its rank, and let \(r_1,r_2,r_3\in G\). For \(i\in\{1,2,3\}\), let \(R_i=\langle r_i\rangle^G\) denote the smallest normal subgroup of \(G\) containing \(r_i\). For subgroups \(A,B\leq G\), write
\[
[A,B]=\langle a^{-1}b^{-1}ab : a\in A,\ b\in B\rangle.
\]
Set \(N=R_1\cap R_2\cap R_3\) and
\[
D_3^{+}=[R_1,R_2\cap R_3]\,[R_2,R_1\cap R_3]\,[R_3,R_1\cap R_2].
\]
These commutator subgroups are normal in \(G\) and contained in \(N\), so their product is a normal subgroup of \(N\).

The source defines, for normal subgroups \(R_1,\ldots,R_n\),
\[
[[R_1,\ldots,R_n]]
=\prod_{\substack{I\cup J=\{1,\ldots,n\}\\ I\cap J=\varnothing}}
\left[\bigcap_{i\in I}R_i,\bigcap_{j\in J}R_j\right].
\]
It does not specify whether \(I\) and \(J\) must both be nonempty. With that restriction, its denominator for \(n=3\) is precisely \(D_3^{+}\).

\emph{Question under the nonempty-partition interpretation.} For every free group \(G\) and every triple \(r_1,r_2,r_3\in G\), is
\[
\frac{R_1\cap R_2\cap R_3}{D_3^{+}}
\]
torsion-free? Explicitly, must
\[
\forall x\in N\ \forall m\in\Z_{\geq1},\qquad
x^m\in D_3^{+}\ \Longrightarrow\ x\in D_3^{+}
\]
hold?

\emph{Unresolved convention.} If empty parts are included and an empty intersection is interpreted as \(G\), the displayed source definition instead gives
\[
D_3^{0}=D_3^{+}[G,N].
\]
The literal question under that convention concerns \(N/D_3^{0}\). The supplied record does not explicitly resolve this distinction; the nonempty-partition interpretation is consistent with its stated two-subgroup comparison.

The source attributes to B. Hartley and Yu. Kuzmin the result that \((R_1\cap R_2)/[R_1,R_2]\) is free abelian when each \(R_i\) is the normal closure of one element of a free group. It also states that the analogous quotient can have nontrivial torsion for \(n\geq4\), and that the three-subgroup question has an affirmative answer when none of \(r_1,r_2,r_3\) is a proper power in \(G\). Here a proper power means an element of the form \(s^k\) with \(s\in G\) and an integer \(k\geq2\). These are supplied background assertions, with the partition convention still requiring confirmation; the question itself imposes no proper-power restriction.
\subsection{Short English statement}
For three elements of any free group, take their normal closures and intersect them. Must the quotient by the product of the three indicated commutator subgroups be torsion-free? The source reports an affirmative answer when none of the elements is a proper power. Its partition notation leaves unresolved whether empty parts contribute additional commutators.
\subsection{Sources}
\begin{itemize}
\item Local UnsolvedMath registry, permanent ID \(701383\), exact source code \texttt{catalogue entry}, title \emph{Kourovka Notebook Problem 17.81}. The supplied record labels the problem \texttt{open}; this is a supplied status, not an independently verified assessment. Upstream revision: \texttt{d371ee917ac9647fbb06cd6fc529deec6f3dda82}; record updated \texttt{2026-10-16T14:00:00Z}.
\item R. Mikhailov, Problem 17.81, \emph{The Kourovka Notebook (Unsolved Problems in Group Theory)}, 17th issue (2010). The supplied record also identifies the 21st issue edition (2026), arXiv identifier \texttt{1401.0300v48}. The collection description names E. I. Khukhro and V. D. Mazurov as editors. These bibliographic details are supplied attribution and have not been independently checked.
\item B. Hartley and Yu. Kuzmin, \emph{J. Pure Appl. Algebra} \textbf{74} (1991), 247--256, cited in the supplied statement for the two-subgroup free-abelian result. No article title or external link is supplied.
\end{itemize}
\end{document}
